\addtocontents{toc}{\protect\Needspace{3\baselineskip}}
\originalchapter{相对论弦的方程与自由端点}\label{ch:basic-NG}
对应 8.251 第 10--11 讲. 目标是从几何面积得到动力学, 并明确哪些速度代表实际形状运动, 哪些只是弦上的重新标记. 此章的张力为 $T_0$, 保留 $c$; 后文自然单位张力 $T_s=1/(2\pi\alpha')$ 与之对应.

\originalsection{张力为何乘世界面面积}
面积的单位是长度平方, 作用量单位是能量乘时间, 因而系数应为 $T_0/c$. 取 Nambu--Goto 作用量
\[
 S_{\rm NG}=-\frac{T_0}{c}\int\dd\tau\dd\sigma\,\mathcal D,
 \quad\mathcal D=\sqrt{(\dot X\cdot X')^2-\dot X^2X'^2}.
\]
若弦静止、长度为 $a$, 面积为 $ca\int\dd t$, 故 $S=-T_0a\int\dd t$, 正好对应静止能量 $E=T_0a$. 这个检验同时固定系数与符号. 洛伦兹不变性来自内积, 二维重参数化不变性来自上一章的行列式和 Jacobian 抵消.

\originalsection{动量密度、方程与边界项}
令 $\mathcal L=-T_0\mathcal D/c$. 对目标坐标导数求偏导, 得
\begin{align*}
 P_\mu^\tau&=\frac{T_0}{c\mathcal D}
 \big[X'^2\dot X_\mu-(\dot X\cdot X')X'_\mu\big],\\
 P_\mu^\sigma&=\frac{T_0}{c\mathcal D}
 \big[\dot X^2X'_\mu-(\dot X\cdot X')\dot X_\mu\big].
\end{align*}
固定初末世界面时间切片后, 积分分部得到
\[
 \delta S=-\int\dd\tau\dd\sigma\,
 (\partial_\tau P_\mu^\tau+\partial_\sigma P_\mu^\sigma)\delta X^\mu
 +\int\dd\tau\,[P_\mu^\sigma\delta X^\mu]_{0}^{\sigma_1}.
\]
内部欧拉--拉格朗日方程为 $\partial_\tau P_\mu^\tau+\partial_\sigma P_\mu^\sigma=0$. 等价的几何式是 $\partial_a(\sqrt{-h}\,h^{ab}\partial_bX^\mu)=0$. 其非线性来自 $h$ 本身也由 $X$ 决定, 不能把 $h$ 当常数后随意解平面波.

开弦边界沿某方向可固定 $X^i$ 并要求 $\delta X^i=0$, 或让其自由变分并要求 $P_i^\sigma=0$. 对通常传播的开弦, 时间方向为 Neumann, 不把端点冻结在一个目标时刻. 闭弦则依靠 $X,P^\sigma$ 和允许变分的周期性使两端贡献抵消, 没有真实的空间端点外力.

\originalsection{D 膜与守恒量的条件}
平直静止 D$p$ 膜可取 $X^{p+1},\ldots,X^{D-1}$ 固定, $X^0,\ldots,X^p$ 为自由方向. $p$ 数的是空间维数, 世界体积维数为 $p+1$. D0 膜是目标空间中的点状对象, 开弦两端即使落在同一点也不因此成为闭弦: 参数拓扑和边界条件仍不同.

总目标动量为 $p_\mu=\int_0^{\sigma_1}P_\mu^\tau\dd\sigma$, 方程给
\[
 \frac{\dd p_\mu}{\dd\tau}=-[P_\mu^\sigma]_{0}^{\sigma_1}.
\]
两端自由的方向动量守恒, 固定方向上膜可施加外力. 若考虑动态膜与弦的联合系统, 总动量应包括膜的反冲; 此时不能只看弦的积分就宣布系统不守恒.

\originalsection{重参数化带来的约束}
由上面的 $P^\tau$ 可直接缩并得
\[
 P^\tau\cdot X'=0,\qquad
 (P^\tau)^2+\left(\frac{T_0}{c}\right)^2X'^2=0.
\]
第一式表示沿弦的标记方向没有独立动量. 对第二式, 展开分子得到
$X'^4\dot X^2-X'^2(\dot X\cdot X')^2=-X'^2\mathcal D^2$, 再除 $\mathcal D^2$ 即可. 约束不是额外选择的经验定律, 而是几何作用量的必然后果. 因拉格朗日量对 $\dot X$ 一次齐次, 朴素规范未固定的正则哈密顿量为零; 真正的目标时间能量由 $-cp_0$ 给出, 不等于“系统没有能量”.

\originalsection{静态规范与横向速度}
取 $X^0=ct$, $\boldsymbol X=\boldsymbol x(t,\sigma)$, 并记弦上物理弧长为 $\ell$. 单位切向量 $\boldsymbol t=\boldsymbol x'/|\boldsymbol x'|$, 横向形状速度为
\[
 \boldsymbol v_\perp=\dot{\boldsymbol x}
 -(\dot{\boldsymbol x}\cdot\boldsymbol t)\boldsymbol t.
\]
随时间重新命名 $\sigma$ 会给 $\dot{\boldsymbol x}$ 增加一个切向分量, 但不会改变 $\boldsymbol v_\perp$. 这不是说普通原子绳子没有纵向运动; 基本无结构弦的参数点没有附加的可观测原子身份, 因而此处的切向标记速度属于规范信息.

代入内积,
\[
 \mathcal D=|\boldsymbol x'|\sqrt{c^2-v_\perp^2},\qquad
 S=-T_0\int\dd t\int\dd\ell\sqrt{1-v_\perp^2/c^2}.
\]
内部类时条件要求 $v_\perp<c$. 静止小段 $\dd\ell$ 的静能为 $T_0\dd\ell$, 对应质量 $T_0\dd\ell/c^2$. 求静态时间的正则动量与能量密度可得
\[
 \frac{\dd\boldsymbol p}{\dd\ell}
 =\frac{T_0}{c^2}\frac{\boldsymbol v_\perp}{\sqrt{1-v_\perp^2/c^2}},
 \qquad\frac{\dd E}{\dd\ell}
 =\frac{T_0}{\sqrt{1-v_\perp^2/c^2}}.
\]
它们说明运动弦段的能量增加, 同时纵向动量保持为零. 参数密度则须乘 $\dd\ell/\dd\sigma$, 不可把每单位 $\sigma$ 当每单位长度.

\originalsection{非相对论横波是面积作用量的展开}
进一步取 $\boldsymbol x=(x,\boldsymbol y(t,x))$, $X^1=x$. 此时
\[
 \frac{\mathcal D^2}{c^2}
 =1+\boldsymbol y_x^2-\frac{\dot{\boldsymbol y}^2}{c^2}
 -\frac{\dot{\boldsymbol y}^2\boldsymbol y_x^2
 -(\dot{\boldsymbol y}\cdot\boldsymbol y_x)^2}{c^2}.
\]
若斜率与 $|\dot{\boldsymbol y}|/c$ 同为小量, 最后一项为四阶. 平方根展开后
\[
 S=-T_0a\int\dd t+
 \frac12\int\dd t\dd x\left[
 \frac{T_0}{c^2}\dot{\boldsymbol y}^2-T_0\boldsymbol y_x^2\right]+O(y^4).
\]
与第\ref{ch:basic-variation}章比较, $\mu_0=T_0/c^2$, 所以波速为 $c$. 同一推导还保留了静能常数 $T_0a$, 避免把小振动能误当整根弦的全部能量.

\originalsection{为什么自由开弦端点是类光的}
在内部可选择共形参数, 使方程为 $\ddot X-X''=0$, 同时保留约束 $\dot X^2+X'^2=0$, $\dot X\cdot X'=0$. 真正自由端点的 Neumann 条件为 $X'=0$. 代入约束得
\[
 \dot X^2=0\qquad\text{在自由端点},
\]
故其世界线类光; 以目标时间衡量, 端点速度大小为 $c$. 这里是约束的端点极限, 没有在 $\det h=0$ 处使用逆度规.

若几何弦有光滑的端点切线, 可从近端点的 $X'$ 方向取极限. 共形条件使端点速度与该切线垂直. 对具有通常偶延拓的 Neumann 解, $X'\simeq\sigma X''$ 而 $\dot X\cdot X'=0$ 的近端点展开给 $\dot X\cdot X''=0$; 这证明的是正则极限的切线垂直性. 若弦在端点形成尖点, 切线本身不唯一, 应只保留类光边界结论.

端点在 D 膜固定方向可承受约束力, 不能把以上“全部方向自由”的结果推广成所有 D 膜端点必须以光速运动. 例如静止绷紧在两块膜之间的弦就有固定端点, 其外力由膜承担.
